\section{Intro}
\label{sec:intro}

In this study, we examine how market fragmentation affects stock liquidity. The theoretical prediction is ambiguous: trading on several venues may improve liquidity through competition for order flow, but it may also reduce depth at each individual venue. We use two regulatory shocks to Swiss-listed equities to study this question. On 1 July 2019, the lapse of EU stock market equivalence forced trading in Swiss shares to consolidate on the SIX Swiss Exchange. On 8 February 2021, UK equivalence was reinstated, partly re-fragmenting the market.

We use SMI constituents as the treatment group and CAC40 and DAX30 constituents as the control group. Our empirical strategy is a difference-in-differences design with stock and day fixed effects. The main outcome variable is the quoted bid-ask spread, while Amihud illiquidity is used as a robustness measure. We first document pre-period differences between the groups, especially lower bid-ask spreads for SMI stocks. These level differences are absorbed by stock fixed effects in the regression analysis.

We find that the July 2019 consolidation widened SMI bid-ask spreads by approximately 2.7 basis points relative to the control group. The effect is statistically significant in our preferred specification with controls and is supported by the Amihud illiquidity results. The partial re-fragmentation in 2021 does not reverse the effect, suggesting that the post-2021 market structure did not restore the liquidity conditions from before the 2019 equivalence lapse.
